% Chapter 19 -- Claim: Most Published Anomalies Fail to Replicate
% Requires: amsmath, amssymb, booktabs, array, tabularx, longtable, graphicx, tcolorbox (breakable, skins), hyperref or url.
% Defines the callout box only if the preamble has not.
\providecolor{fenavy}{HTML}{1B2A6B}
\providecolor{feshade}{HTML}{F3F5FA}
\makeatletter
\@ifundefined{callout}{%
  \newtcolorbox{callout}{breakable, enhanced,
    colback=white, colframe=gray!70, boxrule=0.5pt, arc=1mm,
    borderline west={2.5pt}{0pt}{gray!70},
    left=8pt, right=8pt, top=4pt, bottom=4pt}%
}{}
\makeatother

\providecolor{fenavy}{HTML}{1B2A6B}
\providecolor{feshade}{HTML}{F3F5FA}
\makeatletter
% Shaded box for optional technical material that general readers may skip.
\@ifundefined{technote}{%
  \newtcolorbox{technote}{breakable, enhanced,
    colback=feshade, colframe=fenavy!60, boxrule=0.5pt, arc=1mm,
    left=8pt, right=8pt, top=4pt, bottom=4pt}%
}{}
\makeatother

\chapter{Claim: Most Published Anomalies Fail to Replicate}
\label{ch:16}
\providecommand{\term}[1]{\emph{#1}}

News stories regularly report a new study showing that some simple rule, such as buying cheap stocks or stocks with rising profits, has beaten the market. Other sciences, such as psychology and medicine, have discovered in recent years that many published findings could not be repeated when others tried. This chapter asks whether finance has the same problem, and how a reader can tell a lasting discovery from a lucky one.

By the 2010s academic finance had published hundreds of variables said to predict which stocks would earn more than others. Harvey, Liu and Zhu (2016) counted more than 300. They argued that, because so many had been tried, a new factor should clear a $t$-statistic of about 3 rather than 2. The worry was not new. Chapter~\ref{ch:01} showed how searching produces records that luck alone can explain. What was new was that the worry now applied to a whole literature. McLean and Pontiff (2016) took a different approach. Pure data mining predicts that returns vanish once the original sample ends, while mispricing predicts that they shrink after publication, as investors trade on them. They followed 97 predictors after their original samples ended and after they were published. Average returns were about 26\% lower after the sample ended, and 58\% lower after publication.

Two large replication studies sharpened the debate. Hou, Xue and Zhang (2020) asked how many anomalies survive when tiny stocks count for little, and rebuilt 452 anomalies. They weighted stocks by market value and set group cut-offs using NYSE stocks only, two choices that limit the influence of the tiniest stocks, called micro-caps. About two thirds of the anomalies then failed even the conventional test. Jensen, Kelly and Pedersen (2023) asked whether a verdict that uses the evidence from all factors at once would be kinder, and rebuilt 153 factors in 93 countries. They used a Bayesian method that pulls each factor's estimate toward what similar factors suggest, and concluded that most factors do replicate. Because much of the disagreement lay in each study's code, Chen and Zimmermann (2022) published open code and data for more than 200 predictors, so that anyone could check. Respected researchers using largely the same data reached opposite conclusions, because they asked slightly different questions and applied different standards.

At stake is the credibility of a whole field. If most anomalies are false, much of the factor-investing industry rests on noise. If most are real but fade after publication, investors are learning from research and trading the patterns away. The verdict depends less on the data than on the rules of evidence, so this chapter asks the reader to choose those rules and defend them. The claim on trial\index{claims-on-trial method} is: \textbf{most published cross-sectional return predictors do not replicate once micro-caps, multiple testing and publication are accounted for.}

The lab in Section~\ref{sec:19-lab} uses the free Open Source Asset Pricing data of Chen and Zimmermann. It applies the standards of Hou, Xue and Zhang (2020) and of Jensen, Kelly and Pedersen (2023) to the same predictors, and it measures the decay found by McLean and Pontiff (2016). The goal is to see why careful researchers reached opposite conclusions, and which choices drive the difference.

By the end of the chapter you will be able to:

\begin{itemize}
  \item explain market efficiency\index{efficient market hypothesis}, the joint-hypothesis problem\index{joint-hypothesis problem}, and the kinds of replication, from reproduction\index{reproduction} to pure, scientific and out-of-sample replication\index{replication!pure, scientific and out-of-sample};
  \item explain how publication bias\index{publication bias} inflates published results, and how empirical Bayes\index{empirical Bayes} shrinkage\index{shrinkage} corrects noisy estimates;
  \item set out the prosecution's and the defense's cases, and the design choices that lead them to opposite verdicts;
  \item separate post-publication\index{anomaly!post-publication decay} decay caused by selection from decay caused by investors' learning, and defend your own verdict.
\end{itemize}


\section{Foundations: market efficiency, anomalies and replication}
\label{sec:19-foundations}

\subsection{The efficient market hypothesis}

Section~\ref{sec:01-efficient-markets} introduced the efficient market hypothesis and Fama's (1970) three forms, weak, semi-strong and strong. The debate in this chapter concerns the semi-strong form\index{efficient market hypothesis!semi-strong form}: whether publicly available characteristics of firms, such as their valuation ratios or past returns, predict returns beyond what risk justifies.

\subsection{The joint-hypothesis problem}

Section~\ref{sec:01-efficient-markets} also introduced the joint-hypothesis problem. A strategy that beats a model such as the CAPM\index{capital asset pricing model (CAPM)} shows either that the market is inefficient or that the model is wrong, and the two cannot be separated (Fama 1970). For anomalies this has a direct consequence. One researcher can read the same evidence as mispricing\index{mispricing}, and another as a reward for a risk the model leaves out.

\subsection{Anomalies}

An \textbf{anomaly}\index{anomaly|textbf} is a pattern in returns that a benchmark asset pricing model cannot explain. It is a characteristic, such as a valuation ratio or a past return, that predicts future returns after adjusting for the model's risk factors. Anomalies are usually found with \term{portfolio sorts}\index{portfolio sort} and long--short returns (Section~\ref{sec:05-foundations}). They are judged by the t-statistic\index{t-statistic} of the average return or of the alpha\index{alpha}.

\subsection{Replication, reproduction and robustness}

Terms are used loosely, but a useful distinction is:

\begin{itemize}
  \item \term{Reproduction}\index{reproduction|textbf}: the same data and code give the same numbers. It checks for errors.
  \item \textbf{Pure replication}\index{replication!pure, scientific and out-of-sample|textbf}: the same method on the same kind of data, perhaps rebuilt independently, gives the same result. It checks that the finding is not an artifact of one implementation.
  \item \textbf{Scientific replication} or \term{robustness}\index{robustness|textbf}: different but reasonable choices of method, sample or market give a similar result. It checks that the finding is about the world, not about the choices.
  \item \textbf{Out-of-sample replication}: the result holds in data not used to find it, such as later years or other countries.
\end{itemize}

The two sides of this chapter's debate differ largely in which of these they demand.

\subsection{Publication bias}

Journals are more likely to publish significant results. If only strong in-sample\index{in-sample versus out-of-sample} results are published, the published estimates are too high on average, even if every study was done honestly. The reason is that the published set has been selected for luck as well as for truth. Part E of this chapter's lab measures how large this effect can be.

\subsection{Shrinkage}

Suppose many noisy estimates are made at once. Each one can be improved by pulling it toward the average of all of them. Noisy estimates are pulled more, and precise ones less. A very high estimate is more likely to be partly luck than a moderate one, so pulling it back makes it more accurate on average. This idea comes from \term{empirical Bayes}\index{empirical Bayes|textbf} statistics (Efron and Morris 1975). It is called ``empirical'' because the data themselves supply the average toward which the estimates are pulled. It is the basis of the defense's method.

\textbf{Worked example.} One hundred honest researchers each test a different trading signal, and none of the signals has any real predictive power. By chance, about five will find t-statistics beyond $\pm 1.96$, and those five publish (Figure~\ref{fig:19-pubbias}). Every published paper reports a significant result, yet none of the results is real. A reader who saw only the journals would find a literature with a 100\% success rate.

\begin{figure}[htbp]
\centering
\scalebox{0.9}{%
\begin{tikzpicture}[declare function={npdf(\x) = exp(-\x*\x/2)/sqrt(2*pi);}]
\begin{axis}[
  width=0.68\linewidth, height=0.34\linewidth,
  axis lines=left, axis line style={-stealth},
  xmin=-4, xmax=4.1, ymin=0, ymax=0.66,
  xtick={-1.96,0,1.96}, xticklabels={$-1.96$,$0$,$1.96$},
  ytick=\empty,
  xlabel={t-statistic of a signal with no true power},
  ylabel={Share of signals},
  samples=200, domain=-4:4, clip=false,
]
  \addplot[name path=curve, draw=none] {npdf(x)};
  \path[name path=base] (axis cs:-4,0) -- (axis cs:4,0);
  \addplot[fenavy] fill between[of=curve and base, soft clip={domain=-4:-1.96}];
  \addplot[fenavy] fill between[of=curve and base, soft clip={domain=1.96:4}];
  \addplot[fegray, thick] {npdf(x)};

  \node[align=center] at (axis cs:0,0.56) {about 95 signals: not significant,\\never published};
  \node[fenavy, align=center, font=\footnotesize, anchor=south] (L) at (axis cs:-2.75,0.13) {about 2.5 signals:\\``significant,''\\published};
  \node[fenavy, align=center, font=\footnotesize, anchor=south] (R) at (axis cs:2.75,0.13) {about 2.5 signals:\\``significant,''\\published};
  \draw[-stealth, fenavy] (axis cs:-2.7,0.12) -- (axis cs:-2.3,0.03);
  \draw[-stealth, fenavy] (axis cs:2.7,0.12) -- (axis cs:2.3,0.03);
\end{axis}
\end{tikzpicture}%
}
\caption{If only significant results are published, the journals show only the lucky draws. Publication bias. One hundred signals with no true power produce t-statistics spread around zero. Only the few in the shaded tails, about 2.5 on each side, look significant and get published. The published record shows only the tails.}
\label{fig:19-pubbias}
\end{figure}

\section{How it works: the factor zoo and how replications are designed}
\label{sec:19-how}

\textbf{The zoo.} For decades, finance explained why some stocks earn more than others with a single factor, the market, as in the CAPM. Later it used a handful, such as size, value and momentum\index{momentum}. By 2011 the literature had produced so many candidate predictors that John Cochrane, in his presidential address to the American Finance Association, called it a ``zoo of new factors'' (Cochrane 2011). He argued that the field needed to sort out which factors were truly independent, which were redundant, and which were simply noise. Harvey, Liu and Zhu (2016) counted 316 factors published in leading journals, and the count has kept growing.

The problem is statistical as much as economic. Thousands of researchers test thousands of variables on the same U.S. stock data, so some variables will look significant by luck alone. At the usual cut-off of a t-statistic near 2, about one useless variable in twenty passes, and journals mostly publish the ones that pass (Section~\ref{sec:19-foundations}). The published zoo is therefore a mixture of real effects and false discoveries, much as a stack of winning lottery tickets says little about whether anyone has a winning system. Many of its members are also the same idea measured in slightly different ways, such as several versions of the finding that cheap stocks beat expensive ones. Since Harvey, Liu and Zhu's count, two open data projects have made it possible to rebuild most of the zoo:

\begin{itemize}
  \item Chen and Zimmermann's \textbf{Open Source Asset Pricing}\index{Open Source Asset Pricing|textbf} (2022) replicates more than 200 cross-sectional predictors from the original papers' descriptions and publishes their portfolio returns and code.
  \item Jensen, Kelly and Pedersen (2023) built a \term{global factor database} of 153 factors across 93 countries, with harmonized definitions, meaning each factor is built the same way in every country.
\end{itemize}

No table can list hundreds of predictors, but none needs to, because most of them are variations on a few ideas. Hou, Xue and Zhang\index{q-factor model} (2020) sort their 452 anomalies into six categories. Table~\ref{tab:19-zoo} shows two or three well-known predictors from each. The sign shows whether stocks with a high value of the predictor earned higher (+) or lower ($-$) average returns in the original study. The Open Source Asset Pricing and JKP documentation give the full lists, with definitions and code for every predictor. Jensen, Kelly and Pedersen (2023) group their 153 factors into 13 similar themes.

\begingroup
\small


\renewcommand{\arraystretch}{1.15}
\setlength{\LTpre}{6pt}\setlength{\LTpost}{6pt}
\begin{longtable}{@{}>{\raggedright\arraybackslash}p{0.13\linewidth}>{\raggedright\arraybackslash}p{0.17\linewidth}>{\raggedright\arraybackslash}p{\dimexpr \linewidth-0.13\linewidth-0.17\linewidth-0.21\linewidth-0.05\linewidth-8\tabcolsep\relax}>{\raggedright\arraybackslash}p{0.21\linewidth}>{\centering\arraybackslash}p{0.05\linewidth}@{}}
\caption{A map of the factor zoo: representative predictors in the six categories of Hou, Xue and Zhang (2020).}
\label{tab:19-zoo} \\
\toprule
Category & Predictor & What it measures & Original study & Sign \\
\midrule
\endfirsthead
\multicolumn{5}{@{}l}{\emph{(continued)}} \\
\toprule
Category & Predictor & What it measures & Original study & Sign \\
\midrule
\endhead
\midrule
\multicolumn{5}{r@{}}{\emph{(continued on the next page)}} \\
\endfoot
\bottomrule
\endlastfoot
Momentum & Price momentum & Return over the past 12 months, skipping the most recent month & Jegadeesh and Titman (1993) & $+$ \\
 & Earnings surprise & Change in quarterly earnings from a year earlier, scaled by its volatility & Bernard and Thomas (1989) & $+$ \\
 & Short-term reversal & Return over the past month & Jegadeesh (1990) & $-$ \\
\addlinespace
Value versus growth & Book-to-market & Book value of equity divided by market value & Rosenberg, Reid and Lanstein (1985); Fama and French (1992) & $+$ \\
 & Earnings-to-price & Earnings divided by market value & Basu (1977) & $+$ \\
\addlinespace
Investment & Asset growth & Annual percentage growth in total assets & Cooper, Gulen and Schill (2008) & $-$ \\
 & Net share issuance & Annual change in shares outstanding & Pontiff and Woodgate (2008) & $-$ \\
 & Accruals & Earnings minus operating cash flow, scaled by assets & Sloan (1996) & $-$ \\
\addlinespace
Profitability & Gross profitability & Revenue minus cost of goods sold, divided by assets & Novy-Marx (2013) & $+$ \\
 & F-score & Sum of nine accounting signals of financial strength & Piotroski (2000) & $+$ \\
\addlinespace
Intangibles & R\&D intensity & Research and development spending divided by market value & Chan, Lakonishok and Sougiannis (2001) & $+$ \\
\addlinespace
Trading frictions & Size & Market value of equity & Banz (1981) & $-$ \\
 & Idiosyncratic volatility & Volatility of returns not explained by the Fama--French factors & Ang, Hodrick, Xing and Zhang (2006) & $-$ \\
 & Illiquidity & Average absolute daily return per dollar of trading volume & Amihud (2002) & $+$ \\
 & Maximum daily return & Largest daily return in the past month & Bali, Cakici and Whitelaw (2011) & $-$ \\
\end{longtable}
\endgroup

\textbf{Design choices in a replication\index{replication}.} A replication of an anomaly must make many choices, each defensible, that can change the result:

\begingroup
\small
\setlength{\LTpre}{6pt}\setlength{\LTpost}{6pt}
\begin{longtable}{@{}>{\raggedright\arraybackslash}p{0.17\linewidth}>{\raggedright\arraybackslash}p{\dimexpr 0.9\linewidth-0.17\linewidth-0.34\linewidth-4\tabcolsep\relax}>{\raggedright\arraybackslash}p{0.34\linewidth}@{}}
\toprule
Choice & Options & Why it matters \\
\midrule
\endfirsthead
\multicolumn{3}{@{}l}{\emph{(continued)}} \\
\toprule
Choice & Options & Why it matters \\
\midrule
\endhead
\midrule
\multicolumn{3}{r@{}}{\emph{(continued on the next page)}} \\
\endfoot
\bottomrule
\endlastfoot
Universe & All stocks, or excluding micro-caps (for example below the NYSE 20th size percentile) & Many anomalies are concentrated in tiny stocks that are costly to trade and hold little value \\*
Weighting & Equal-weighted or value-weighted portfolios & Equal weighting gives micro-caps large influence \\
Breakpoints & Deciles, quintiles or terciles; NYSE or all-stock breakpoints & Extreme portfolios magnify signals and noise \\
Sample & The original paper's sample, or the longest available & Later years are out of sample \\
Significance hurdle & $t > 1.96$, $t > 3$, a false-discovery-rate rule, or a Bayesian posterior & Determines what counts as a success \\*
Benchmark & Raw returns, CAPM alpha, or alpha against a multifactor model & A factor explained by others may be real but redundant \\
\end{longtable}
\addtocounter{table}{-1}% uncaptioned longtable: keep table numbers unchanged
\endgroup

\term{NYSE breakpoints}\index{NYSE breakpoints|textbf} set the cut-offs between groups using only stocks listed on the New York Stock Exchange. These stocks are larger on average, so the many tiny Nasdaq stocks do not crowd the extreme portfolios. A \term{Bayesian posterior}\index{Bayesian posterior|textbf} is the probability of a hypothesis after a prior belief has been combined with the data.

\textbf{Pure and scientific replication again.} A replication that copies the original method tests whether the paper's numbers can be reproduced. One that changes the universe, weighting or hurdle tests whether the finding is robust. Both are valuable, but they answer different questions, and the debate in this chapter turns partly on that difference.

\section{The claim}
\label{sec:19-claim}

Stated so that it can fail:

\begin{callout}
\textbf{Most published cross-sectional return predictors do not reflect a real, persistent relationship: they fail to replicate under reasonable standards, or their returns disappear outside the sample in which they were found.}
\end{callout}

``Reasonable standards'' is where the debate lies. The claim can be tested by \textbf{replication} within the original sample (which hurdle, and which portfolio construction?), \textbf{out of sample}\index{in-sample versus out-of-sample|textbf} (do returns persist after the sample ends and after publication?), and \term{statistical correction} (after allowing for the number of predictors tested and for publication bias, what share probably have a positive true return?).

\section{The prosecution}
\label{sec:19-prosecution}

\textbf{Harvey, Liu and Zhu (2016)} argued that because hundreds of factors have been tested, and many more never published, the usual hurdle of $t > 2$ lets through too many false discoveries. They proposed a hurdle of about 3 and concluded that many published factors were likely false.

\textbf{Hou, Xue and Zhang (2020)}\index{q-factor model|textbf} rebuilt 452 anomalies with one method designed to limit the influence of micro-caps\index{micro-caps}: NYSE breakpoints and value-weighted\index{index weighting!value-weighted} returns. 65\% failed a t-statistic of 1.96, and 82\% failed a multiple-testing hurdle of 2.78. Anomalies based on \term{trading frictions}\index{trading frictions|textbf}, such as illiquidity\index{illiquidity} and \term{idiosyncratic risk}\index{risk!idiosyncratic|textbf} (the part of a stock's volatility\index{volatility} that the market and other factors do not explain), failed most often, 96\% of them. They concluded that most anomalies are false or confined to micro-caps too small and costly to matter.

The prosecution rests on two premises: the hurdle should be high because many tests have been run, and portfolios should give little weight to the smallest stocks.

\section{The defense}
\label{sec:19-defense}

\textbf{Jensen, Kelly and Pedersen (2023)} used a Bayesian model that shrinks each factor's noisy alpha toward the average of related factors, rather than judging it alone against a fixed hurdle. They found that most factors replicate, cluster into 13 themes that are mostly significant parts of the optimal portfolio, work out of sample in 93 countries, and are corroborated rather than weakened by their large number.

\textbf{Chen and Zimmermann (2022)} found that almost all predictors whose original papers reported clear evidence can be reproduced with the original methods. \textbf{Chen and Zimmermann (2020)} estimated that, across 156 predictors, publication bias makes in-sample returns only about 12\% too large, and that a t-statistic hurdle of about 1.8 is enough for published predictors. \textbf{Chen (2021)} argued that published t-statistics are too large to be explained by \term{p-hacking}\index{p-hacking|textbf}, trying analyses until one crosses the significance line, alone.

The defense rests on two different premises: the question is whether the original findings are real, not whether they survive a method that plays down small stocks, and evidence should be pooled across related factors.

\section{Post-publication decay}
\label{sec:19-decay}

\textbf{McLean and Pontiff (2016)}\index{anomaly!post-publication decay|textbf} followed 97 predictors and found returns 26\% lower after the original sample ended and 58\% lower after publication (Figure~\ref{fig:19-decay}). Nothing changed when a sample ended except that researchers stopped looking, so the first figure is an upper bound on data mining\index{data mining}. The further decline after publication reflects investors trading on the results, which suggests the original findings were real. Both forces matter, and the publication effect is the larger.

\begin{figure}[htbp]
\centering
\scalebox{0.9}{%
\begin{tikzpicture}
\begin{axis}[
  width=0.8\linewidth, height=0.34\linewidth,
  ybar, bar width=0.68cm,
  axis lines=left, axis line style={-stealth},
  ymin=0, ymax=118, ytick={0,50,100},
  ylabel={Return (in sample $= 100$)},
  symbolic x coords={A,B,C},
  xtick=data,
  xticklabels={In sample, {After sample,\\before\\publication}, {After\\publication}},
  xticklabel style={align=center, text width=2.3cm, font=\small},
  enlarge x limits=0.25,
  nodes near coords,
]
  \addplot[fill=fenavy, draw=none] coordinates {(A,100) (B,74) (C,42)};
\end{axis}
\end{tikzpicture}%
}
\caption{Published patterns lost more than half of their returns after publication. Post-publication decay in McLean and Pontiff (2016). With in-sample returns set to 100, the average predictor earned 74 between the end of its original sample and publication, and 42 after publication.}
\label{fig:19-decay}
\end{figure}

The lab adds a check the original study could not make: a simulated literature in which the truth is known. If only significant results are published, returns appear to fall after the sample ends even when nothing changes, which gives a benchmark for how much of the observed decay is selection.

\section{Why the two sides disagree}
\label{sec:19-why}

The two sides agree on most facts. They disagree on four choices.

\begingroup
\small
\setlength{\LTpre}{6pt}\setlength{\LTpost}{6pt}
\begin{longtable}{@{}>{\raggedright\arraybackslash}p{0.17\linewidth}>{\raggedright\arraybackslash}p{\dimexpr(0.9\linewidth-0.17\linewidth-4\tabcolsep)/2\relax}>{\raggedright\arraybackslash}p{\dimexpr(0.9\linewidth-0.17\linewidth-4\tabcolsep)/2\relax}@{}}
\toprule
Choice & Prosecution (Hou, Xue and Zhang; Harvey, Liu and Zhu) & Defense (Jensen, Kelly and Pedersen; Chen and Zimmermann) \\
\midrule
\endfirsthead
\multicolumn{3}{@{}l}{\emph{(continued)}} \\
\toprule
Choice & Prosecution (Hou, Xue and Zhang; Harvey, Liu and Zhu) & Defense (Jensen, Kelly and Pedersen; Chen and Zimmermann) \\
\midrule
\endhead
\midrule
\multicolumn{3}{r@{}}{\emph{(continued on the next page)}} \\
\endfoot
\bottomrule
\endlastfoot
Which stocks count & De-emphasize micro-caps: NYSE breakpoints, value weights & Replicate the original construction; micro-caps are part of the market \\*
Hurdle & High: $t$ of about 3 or 2.78, for multiple testing & Lower, or a Bayesian posterior \\
Unit of evidence & Each factor judged on its own & Related factors pooled; corroboration strengthens the evidence \\*
Kind of replication & Scientific: does it survive reasonable alternative choices? & Pure: are the original findings reproducible and real? \\
\end{longtable}
\addtocounter{table}{-1}% uncaptioned longtable: keep table numbers unchanged
\endgroup

Neither side is wrong on its own terms. A factor found only among micro-caps may be real but irrelevant to an investor who cannot trade them; pooling related factors is more accurate if its assumptions hold, but misleading if a whole theme was mined together. ``Does this anomaly replicate?'' is not well posed until the purpose is stated: pure replication and pooled evidence suit the question of how markets price risk, while value-weighted, liquid, out-of-sample, after-cost evidence suits the question of what to invest in.

\section{Lab~\thechapter: Replicating both sides of the anomaly debate}
\label{sec:19-lab}

The lab replicates Hou, Xue and Zhang (2020), Jensen, Kelly and Pedersen (2023) and McLean and Pontiff (2016). It asks what share of published predictors survive when each side's standards are applied to the same data, and how much of their return remains after their samples end and after publication. The data are the Open Source Asset Pricing monthly long--short returns for more than 200 predictors as constructed in the original papers, with each paper's sample period and publication year, and the code is in the repository folder \texttt{ch16-anomalies/}.

\begin{itemize}
  \item \textbf{Part A, replication.} Each predictor's t-statistic over its original sample, and the share clearing $t > 1.96$, $t > 3$\index{t-statistic!$t>3$ hurdle} and a Benjamini--Hochberg test.
  \item \textbf{Part B, decay.} Average returns in sample, after the sample and after publication, with a \textbf{pooled regression}\index{regression!pooled|textbf} across all predictors in the style of McLean and Pontiff, with errors clustered by\index{standard error!clustered} month.
  \item \textbf{Part C, methodology.} The same statistics for an alternative construction, if supplied, the core of the disagreement.
  \item \textbf{Part D, shrinkage.} Empirical Bayes estimates of each predictor's true mean, and whether shrunk estimates predict post-publication returns better than raw ones.
  \item \textbf{Part E, known truth.} A simulated literature that publishes only signals with in-sample $t > 2$. With \textbf{no} true change after publication, it still shows returns about 54\% lower after publication, purely from publishing lucky results.
\end{itemize}

The lab ends by asking you to reach your own verdict on the scale of Section~\ref{sec:01-method}, state the standard of evidence you apply, and compare your verdict with Section~\ref{sec:19-verdict}.

\section{Where the evidence stands}
\label{sec:19-verdict}

The published studies point to a split picture: \textbf{most published anomalies are not pure statistical flukes. But most are smaller than reported, concentrated in small stocks, partly traded away, and not profitable to trade after costs.} The table summarizes the evidence as of November 2026.

{\small
\begin{longtable}{@{}>{\raggedright\arraybackslash}p{0.19\linewidth}>{\raggedright\arraybackslash}p{0.13\linewidth}>{\raggedright\arraybackslash}p{\dimexpr 0.45\linewidth-6\tabcolsep\relax}>{\raggedright\arraybackslash}p{0.23\linewidth}@{}}
\toprule
Situation & Holds? & Why & Evidence \\
\midrule
\endfirsthead
\multicolumn{4}{@{}l}{\emph{(continued)}} \\
\toprule
Situation & Holds? & Why & Evidence \\
\midrule
\endhead
\midrule
\multicolumn{4}{r@{}}{\emph{(continued on the next page)}} \\
\endfoot
\bottomrule
\endlastfoot
Original construction, conventional hurdle & \textbf{No} & Almost all predictors with clear original evidence reproduce with the original methods. & Chen and Zimmermann (2022) \\*
\addlinespace[3pt]
Micro-caps de-emphasized, value-weighted & \textbf{Yes} & 65\% of 452 anomalies fail $t = 1.96$ and 82\% fail 2.78; 96\% of trading-frictions anomalies fail. & Hou, Xue and Zhang (2020) \\
\addlinespace[3pt]
The right hurdle for multiple testing & \textbf{Disputed} & About 3 on one side; about 1.8 for published predictors on the other. & Harvey, Liu and Zhu (2016); Chen and Zimmermann (2020) \\
\addlinespace[3pt]
Bayesian pooling across factors and countries & \textbf{No} & Most factors replicate, cluster into 13 themes and work out of sample in 93 countries. & Jensen, Kelly and Pedersen (2023) \\
\addlinespace[3pt]
Out of sample, after publication & \textbf{Partly} & Returns were 26\% lower after the samples and 58\% lower after publication, partly from selection. & McLean and Pontiff (2016); Part E of the lab \\*
\addlinespace[3pt]
After costs, for an investor & \textbf{Yes} & Little expected return remains for the average anomaly. & Chen and Velikov (2023); Chapter~\ref{ch:14} \\
\end{longtable}
\addtocounter{table}{-1}% uncaptioned longtable: keep table numbers unchanged
}

Whether this makes most anomalies ``false'' depends on whether the question is scientific or practical, so a verdict on this claim is a choice of standard as much as a reading of the data (Section~\ref{sec:19-why}).

\textbf{What would change the picture.} A shared standard, such as pre-registered\index{pre-registration} hurdles fixed before the data are examined, or value-weighted, micro-cap-robust results reported alongside the original construction, would narrow the disagreement more than any new data set.

\label{sec:19-open}
\textbf{Open questions.} How many signals were tested but never published? The unknown size of this \textbf{file drawer}\index{publication bias|textbf} limits every multiple-testing correction. And is post-publication decay learning or selection, a question Part E answers in simulation but not yet in real data?

\section{What practitioners and investors can use}
\label{sec:19-investors}

To judge any new factor fund or signal, check whether the original t-statistic clears about 3, whether the effect holds among large, liquid stocks with value weights, and whether it belongs to a theme that works, such as value, momentum, profitability or low risk. Expect roughly half the backtest\index{backtest} after publication at best, ask for turnover\index{turnover} and net returns at realistic costs (Chapter~\ref{ch:14}), and ask how many signals were tested before this one.

\begin{callout}
\textbf{What Failed: the trading-frictions anomalies.} Anomalies based on liquidity\index{liquidity}, idiosyncratic volatility\index{idiosyncratic volatility} and other trading frictions looked strong in the papers that introduced them. When Hou, Xue and Zhang played down micro-caps with NYSE breakpoints and value weights, 96\% of them failed to clear a t-statistic of 1.96, because their returns had come mostly from the smallest, least liquid stocks, where costs are largest. The claim that frictions-based anomalies offer returns an investor can capture \textbf{fails}.
\end{callout}

% Key terms of the chapter, in alphabetical order, for readers new to finance.
% A term defined in an earlier chapter points to the chapter, section and page.
\providecolor{FEMaroon}{RGB}{128,30,45}
\providecolor{FECream}{RGB}{253,248,238}
\providecolor{FEGold}{RGB}{196,150,60}
\makeatletter
\@ifundefined{keyterms}{%
  \newtcolorbox{keyterms}[1][Key terms in this chapter]{breakable, enhanced,
    colback=FECream, colframe=FEGold, boxrule=0.6pt, arc=2pt, leftrule=3pt,
    left=8pt, right=8pt, top=6pt, bottom=6pt, before skip=12pt, after skip=12pt,
    fonttitle=\bfseries, coltitle=white, colbacktitle=FEMaroon, title={#1}}%
}{}
\makeatother

\begin{keyterms}
\small
\begin{description}[leftmargin=0pt, itemsep=2pt, parsep=0pt, font=\normalfont\bfseries\color{FEMaroon}]
  \item[Bayesian posterior.] A Bayesian posterior is the probability of a hypothesis after a prior belief has been combined with the data.
  \item[Empirical Bayes.] Empirical Bayes pulls each estimate toward the average of all the estimates, by an amount the data themselves determine, which makes the estimates more accurate on average.
  \item[Global factor database.] The global factor database of Jensen, Kelly and Pedersen (2023) holds 153 factors across 93 countries, each built the same way everywhere.
  \item[Idiosyncratic risk.] Called idiosyncratic variance there. See Chapter~\ref{ch:03}, Section~\ref{sec:04-capm}, page~\pageref{kt:idiosyncratic-variance}.
  \item[NYSE breakpoints.] NYSE breakpoints set the cut-offs between groups in a portfolio sort using only stocks listed on the New York Stock Exchange, so that small stocks do not crowd the extreme groups.
  \item[p-hacking.] p-hacking means trying analyses until one crosses the significance line.
  \item[Portfolio sort.] See Chapter~\ref{ch:04}, Section~\ref{sec:05-foundations}, page~\pageref{sec:05-foundations}.
  \item[Reproduction and replication.] Reproduction checks that the same data and code give the same numbers, while pure and scientific replication check the result with new data or reasonable new methods.
  \item[Robustness.] A result is robust when different but reasonable choices of method, sample or market give a similar answer.
  \item[Statistical correction.] A statistical correction adjusts results for the number of predictors tested and for publication bias.
  \item[Trading frictions.] Trading frictions are the costs and limits that make trading difficult, such as illiquidity, spreads and shorting costs.
\end{description}
\end{keyterms}
